\section{Results}
\subsection{Reporting rebounded partially across hurricane-related periods}
Of \nChecklists{} checklists, \nDetections{} reported the species, an overall rate of \nEncounterPct{}\%. Reporting declined from 7.59\% before Mar\'ia to 4.62\% between hurricanes, then increased to 6.21\% after Fiona (Table~\ref{tab:periods}; Figure~\ref{fig:rates}). The later increase therefore did not restore the aggregate initial rate. Annual summaries show a low of 3.80\% in 2019 and 6.62\% in 2025. The period-level rate recovered by 1.59 percentage points but remained 1.38 points below the initial value, indicating a partial rebound following the decline. This aggregate change raises the environmental question of whether rebound is widespread across previously reporting sites or concentrated in particular elevation and greenness combinations. The stratified comparisons below locate the pattern before site-transition results examine its persistence interpretation.
\inctab{periods_rewrite}
\begin{figure}[htbp]\centering
\incfig{fig03_encounter_rates}
\caption{Observed annual and period-level reporting rates from the released records. Fractions printed above the period bars give positive reports divided by all checklists. Annual intervals are the original binomial intervals and do not account for dependence among checklists.}\label{fig:rates}
\end{figure}

\subsection{\ifzh NDVI 缺失集中于高报告率的山地环境\else NDVI missingness concentrates in high-reporting upland environments\fi}
\ifzh
NDVI 可用比例从 María 前的 72.58\% 增至 Fiona 后的 79.40\%。完整记录的正报告率为 4.89\%（3,370/68,976），缺失记录为 8.71\%（1,804/20,706）。不同海拔层的缺失比例存在明显差异（表~\ref{tab:envmissing}；图~\ref{fig:envmissing}）。低、中、高海拔的缺失率分别为 20.8\%、32.8\% 和 49.6\%。高海拔清单占缺失子集的 11.2\%，但仅占完整子集的 3.4\%；与此同时，该层完整记录的报告率达到 37.4\%，远高于低海拔的 3.3\%。因此，缺失子集的较高总体报告率与高报告率山地记录的较大权重密切相关。
\else
NDVI availability increased from 72.58\% before Mar\'ia to 79.40\% after Fiona. Reporting was 4.89\% among complete records (3,370/68,976) and 8.71\% among missing records (1,804/20,706). Missingness differed considerably among elevation bands (Table~\ref{tab:envmissing}; Figure~\ref{fig:envmissing}). It was 20.8\%, 32.8\% and 49.6\% in low, middle and high bands. High-elevation checklists formed 11.2\% of missing records but only 3.4\% of complete records. Their complete-record report rate was 37.4\%, compared with 3.3\% at low elevation. High-reporting upland environments consequently receive greater weight in the missing subset.
\fi
\inctab{environment_missingness}
\ifzh
这种组成解释也得到条件比较支持。缺失指示变量的未校正报告优势比为 1.86（95\% 区间 1.51--2.28），加入海拔与时期后为 1.03（0.85--1.26）；继续控制努力、季节和位置后为 1.20（1.03--1.41）。较大的总体差异因而主要体现环境与时期组成，进一步校正后仍保留较小的观测差异。在高海拔层内，缺失记录报告率实际略低于完整记录（35.7\% 对 37.4\%），说明缺失并不对应所有环境中一致增加的报告机会。对“哪里回升更明显”的判断而言，完整记录减少了高报告率山地环境的权重，而 NDVI 可用比例又随时期变化；原始汇总因而混合了地点报告变化与环境代表性变化。分层并统一协变量分布，有助于提高时期比较的可比性；补充山地缺失记录的环境信息，则有助于检验较高持续性关联能否覆盖当前环境分析遗漏的清单。
\else
Conditional comparisons support this compositional explanation. The unadjusted report odds ratio for missing records was 1.86 (95\% interval 1.51--2.28), declining to 1.03 (0.85--1.26) after elevation and period adjustment. Further adjustment for effort, season and location gave 1.20 (1.03--1.41). Much of the aggregate difference therefore reflects environmental and period composition, with a smaller residual observation difference after fuller adjustment. Within the high-elevation band, missing records actually had slightly lower reporting than complete records (35.7\% versus 37.4\%), showing that missingness does not correspond to uniformly greater reporting across environments. For locating rebound, complete records downweight high-reporting uplands while NDVI availability also changes across periods. Raw summaries consequently mix changes in reporting with changes in environmental representation. Stratification and a common covariate distribution improve period comparability; recovering environmental information for missing upland records would help assess whether the higher-persistence association extends to observations omitted from the environmental analysis.
\fi
\begin{figure}[htbp]\centering
\incfig{fig25_environment_missingness}
\caption{\ifzh
NDVI 缺失的环境结构。（a）各海拔层缺失比例；（b）层内完整与缺失记录的报告率；（c）逐步校正的报告优势比及空间区块聚类 95\% 区间。
\else
Environmental structure of NDVI missingness. Panels show (a) missing fraction by elevation, (b) reporting within each band, and (c) sequentially adjusted odds ratios with spatial block-clustered 95\% intervals.
\fi}\label{fig:envmissing}
\end{figure}
\ifzh
对占域持续性研究而言，完整记录筛选减少的是山地环境中的清单权重，而不只是均匀减少样本量。占域网格核查显示，三个时期都有记录的 1,048 个网格中只有四个完全没有 NDVI；动态模型仍覆盖 1,044 个网格，且采用地点跨时期平均绿度。缺失数据对清单模型与占域模型的影响有所不同。完整记录筛选改变了清单模型中的山地记录比例，占域模型则更多依赖跨时期平均绿度能否代表各时期的地点条件。监测应保留山地重复访问，同时改善鸟类调查与同期植被观测的匹配，避免把完整清单中的环境权重直接视为全体地点的权重。
\else
For persistence research, complete-case filtering reduces the checklist weight of upland environments rather than removing an environmentally uniform sample. Of 1,048 grids represented in all three periods, only four lacked NDVI entirely; the dynamic model retained 1,044 grids using cross-period mean greenness. Missingness affects the checklist and occupancy analyses differently. Complete-case filtering changes upland record representation in the checklist model, whereas the occupancy model depends on how well cross-period mean vegetation represents conditions in each period. Monitoring should retain repeated upland visits and improve matching to contemporaneous vegetation observations, rather than treating complete-checklist weights as the weights of all sites.
\fi
\FloatBarrier


\subsection{The period contrast has an elevation structure}
The reproduced interaction GLM m4 had AIC \aicMfour{}, model weight \wMfour{} and pseudo-$R^2$ \pseudoRMfour{} (Table~\ref{tab:glmsel}). The additive environmental model m3 remained competitive ($\Delta$AIC 2.49), whereas adding environment to the period-and-effort model reduced AIC by 3704.25. Environmental gradients account for substantial fitted variation within this candidate set; the additional evidence for the period--elevation interaction is more limited. Environmental baselines are therefore central to explaining report distributions, whereas locating stronger rebound requires within-environment period contrasts. Persistently high reporting at greater elevation does not itself establish faster recovery.
\inctab{glmsel}

At fixed covariates, reporting probability at 100~m was 3.77\%, 2.40\% and 3.08\% across the three periods. At 900~m it was 40.58\%, 36.52\% and 44.89\%. Post-Fiona predictions therefore remained below the initial baseline at low elevation while exceeding it at high elevation (Figure~\ref{fig:gradients}). These conditional comparisons and the aggregate incomplete rebound describe different quantities and can hold simultaneously. The odds ratios per 100~m were 1.430, 1.483 and 1.500. With the original model-based standard errors, the Post-Fiona interaction had $p=0.011$; the Inter-Hurricanes interaction had $p=0.063$.

NDVI was positively associated with reporting ($\beta=0.871$, SE 0.087). An increase of 0.1 NDVI corresponded to an odds ratio of approximately 1.091 (95\% Wald interval 1.073--1.110). The probability-weighted elevation centroid shifted by +28.53~m between hurricanes and +13.88~m after Fiona relative to the initial period. These shifts summarize reweighting of the fitted elevation curve, not individual displacement. The within-Post-Fiona linear time term did not reach the 0.05 level ($p=0.121$). Period means and annual trends describe different temporal scales; rebound here denotes the increase between period-level reporting rates.
\begin{figure}[htbp]\centering
\incfig{fig21_reporting_gradients_rewrite}
\caption{Environmental structure of the reproduced reporting pattern. (a) GLM predictions with NDVI, duration and distance fixed at their complete-sample means. (b) Probability-weighted centroid shifts on the uniform 0--1300~m prediction grid. Curves are fitted reporting responses; centroid shifts do not measure movement.}\label{fig:gradients}
\end{figure}
\FloatBarrier

\subsection{\ifzh 环境分层将报告回升定位于不同起点和变化幅度\else Environmental strata locate rebound within different baselines and changes\fi}
\ifzh
九个环境组合的标准化预测显示，报告概率的环境差异远大于同一组合内的时期变化（表~\ref{tab:envrebound}；图~\ref{fig:envrebound}）。低海拔低绿度组合的三个时期概率为 2.29\%、1.42\% 和 1.81\%，低海拔高绿度为 6.27\%、4.10\% 和 5.17\%；两者均表现为下降后部分回升。高海拔高绿度组合则为 42.48\%、38.01\% 和 44.98\%。在相同海拔层内，较高绿度对应更高报告基线；跨海拔比较时，山地组合从较高基线上维持较高报告水平。从两次飓风之间到 Fiona 后，低海拔低、高绿度组合分别增加 0.39 和 1.07 个百分点，高海拔高绿度组合增加 6.97 个百分点。因此，以绝对报告概率增量衡量，较明显的回升点估计出现在高海拔高绿度环境；低地高绿度环境虽有较高报告水平，其末期预测仍低于自身初期。后续分析据此检验山地高报告是否对应较高的占域持续概率，以及低地较绿地点是否在整体报告回升不足的背景下仍具有较高的占域持续概率。
\else
Standardized predictions across the nine environmental cells show larger environmental contrasts than within-cell period changes (Table~\ref{tab:envrebound}; Figure~\ref{fig:envrebound}). Low-elevation, low-greenness probabilities were 2.29\%, 1.42\% and 1.81\% across periods, versus 6.27\%, 4.10\% and 5.17\% at low elevation and high greenness; both partially rebounded after declining. The high-elevation, high-greenness cell gave 42.48\%, 38.01\% and 44.98\%. Higher greenness corresponded to a higher reporting baseline within elevation bands, while upland cells retained high reporting from an already high baseline. From Inter-Hurricanes to Post-Fiona, the low-elevation low- and high-greenness cells increased by 0.39 and 1.07 percentage points, respectively, compared with 6.97 points in the high-elevation, high-greenness cell. By absolute probability change, the larger rebound point estimate therefore occurs in high, green environments. Greener lowlands retain higher reporting levels but remain below their own initial prediction. The subsequent persistence analysis examines whether high upland reporting corresponds to higher occupancy persistence probabilities and whether greener lowlands also have higher persistence despite incomplete aggregate reporting rebound.
\fi
\inctab{environment_rebound}
\ifzh
空间聚类区间显示跨期增量的估计精度较低。低海拔高绿度组合的 Fiona 后相对初期差值为 $-1.10$ 个百分点，95\% 区间为 $[-2.62,0.42]$；高海拔高绿度为 $+2.50$ 个百分点，区间为 $[-5.09,10.09]$。高海拔低绿度仅有 144 条完整清单，而低海拔低绿度有 22,433 条，环境覆盖高度不均衡。当前数据因此更清楚地定位了高、低报告环境，对各环境层已经超越初期水平的估计则较宽。所有环境组合的末期相对初期区间均跨越零，削弱了“山地已恢复而低地尚未恢复”的确定性划分；较高报告基线也会影响绝对增量大小。监测应通过同一地点、相近努力的重复调查，同时追踪各层的绝对变化和相对基线变化，并补充山地低绿度等稀疏组合，检验当前梯度是否反映更广泛的地点变化。
\else
Spatially clustered intervals show lower precision for the temporal increments. The low-elevation, high-greenness Post-Fiona contrast was $-1.10$ percentage points (95\% interval $[-2.62,0.42]$), while the high-elevation, high-greenness contrast was $+2.50$ points ($[-5.09,10.09]$). Coverage was uneven. Only 144 complete checklists represented high elevation with low greenness, versus 22,433 at low elevation with low greenness. The data locate high- and low-reporting environments more precisely than they establish how far each has surpassed its initial level. All cell-specific Post-Fiona versus Pre-Mar\'ia intervals span zero, weakening a categorical claim that uplands have recovered while lowlands have not. Higher baselines also affect the size of absolute increments. Repeated surveys at the same sites with comparable effort should track both absolute and baseline-relative changes and improve coverage of sparse cells, including less-green uplands, to test how widely the fitted gradient extends across sites.
\fi
\begin{figure}[htbp]\centering
\incfig{fig26_environment_rebound}
\caption{\ifzh
环境组合的标准化报告概率及变化。前两图为时期概率，第三图为 Fiona 后减去 María 前的百分点差及空间聚类 95\% 区间。所有时期采用相同的层内协变量分布。
\else
Standardized reporting probabilities and changes by environmental cell. The first two panels show period probabilities; the third gives Post-Fiona minus Pre-Mar\'ia percentage points and spatially clustered 95\% intervals. All periods share the same within-cell covariate distribution.
\fi}\label{fig:envrebound}
\end{figure}
\FloatBarrier


\subsection{Persistence associations remain after observation-model correction}
The ordering refit converged and changed AIC from \aicOccPub{} to \aicOccFix{}. Realigning period labels to the visits at each site changed the fitted observation process. Within the corrected model, initial occupancy remained positively associated with elevation ($\beta=1.647$) and NDVI ($\beta=0.583$), while local extinction remained negatively associated with elevation ($\beta=-0.720$) and NDVI ($\beta=-0.818$). Their approximate 95\% intervals excluded zero (Table~\ref{tab:occintervals}). Thus the environmental association with persistence survives this specific implementation correction.
\inctab{occupancy_intervals}
\inctab{occ}

The detection contrast changed more substantially. The Pre-Mar\'ia coefficient relative to Inter-Hurricanes changed from $-0.566$ (SE 0.190, $p=0.0029$) to $-0.055$ (SE 0.159, $p=0.727$). Corrected detection probabilities were 30.3\%, 31.5\% and 33.8\% across the three periods, with neither period contrast clearly separated from zero (Figure~\ref{fig:ordering}). The absolute initial-period coefficient shrank by approximately 90\%, and the corrected probabilities spanned 3.5 percentage points. The initial detection contrast was consequently sensitive to label ordering, while environmental persistence associations remained stable. Full published, reproduced and corrected coefficients are retained in Table~\ref{tab:occ}.

The NDVI colonization coefficient was close to zero ($\beta=0.057$, 95\% interval $[-0.278,0.392]$). Evidence linking greenness to occupancy persistence is therefore clearer than evidence linking it to entry into unoccupied sites. This strengthens the interpretation linking high reporting in greener environments to retention of occupied sites, while providing less support for attributing rebound to greenness-associated colonization. Within this model, greenness primarily characterized retention of occupied sites, whereas a single greenness gradient explained entry into occupancy weakly. Monitoring should retain detections and nondetections from each visit to estimate occupancy persistence and colonization probabilities separately.
\begin{figure}[htbp]\centering
\incfig{fig16_occupancy_obscovs_ordering}
\caption{Effect of correcting period-label alignment in the selected dynamic occupancy model. The archived coefficient plot uses one-SE error bars; Table~\ref{tab:occintervals} gives 95\% intervals. Detection probabilities follow the panel's displayed period order. Environmental and detection coefficients differ in their sensitivity to observation-order correction.}\label{fig:ordering}
\end{figure}
\FloatBarrier

\subsection{\ifzh 低海拔网格的占域持续概率随绿度而异\else Low-elevation occupancy persistence varies with greenness\fi}
\ifzh
将校正后的转移概率投射到实际地点环境，得到比单独报告系数更明确的持续性结构（表~\ref{tab:envpersistence}；图~\ref{fig:envpersistence}）。低海拔低、中、高绿度地点的平均占域持续概率分别为 24.2\%、49.7\% 和 71.3\%，对应 320、430 和 158 个地点；高海拔高绿度的 39 个地点为 96.8\%。这说明持续性梯度不仅体现在山地与低地之间，低地内部也存在明显的绿度相关分化。与报告分层结果对照，高海拔高绿度环境同时具有较高报告水平和模型估计的较高占域持续概率，支持在这些环境中设置固定调查网格，检验后续时期的占域状态。在低海拔高绿度环境中，报告概率的点估计尚未回到初期水平，但地点占域持续概率相对较高。报告回升程度与已占据地点的保留描述不同过程，低地整体报告偏低不能直接推导为所有低地地点持续性较弱。这些分层概率由同一校正模型推导，提供的是环境关联的具体表达。
\else
Projecting corrected transition probabilities onto observed site environments reveals a persistence structure beyond the coefficient signs (Table~\ref{tab:envpersistence}; Figure~\ref{fig:envpersistence}). Mean occupancy persistence probability was 24.2\%, 49.7\% and 71.3\% across low-elevation low-, middle- and high-greenness cells, containing 320, 430 and 158 sites. The 39 high-elevation, high-greenness sites averaged 96.8\%. Persistence therefore varies both between uplands and lowlands and along greenness within lowlands. Compared with the reporting strata, high-elevation, high-greenness environments combine high reporting with high modeled persistence, identifying environments in which fixed survey grids can test occupancy in later periods. In greener lowlands, the reporting point estimate remains below its initial baseline despite relatively high occupancy persistence probability. Rebound and retention of occupied sites describe different processes, so low aggregate reporting does not imply uniformly weak persistence across lowlands. These probabilities express environmental associations derived from the same corrected model.
\fi
\inctab{environment_persistence}
\ifzh
固定海拔在 100~m 后，NDVI 从 0.30 增至 0.76 时，占域持续概率预测由 29.8\% 升至 72.0\%，差值为 42.2 个百分点（95\% 区间 9.9--64.9）；定殖预测仅由 5.3\% 变为 6.0\%，差值 0.7 个百分点，区间为 $[-2.9,5.3]$。该对照保持海拔不变，支持绿度主要对应已占据地点保留的解释。对保护管理而言，维持已有占域地点的植被条件，与促进未占据地点进入占域状态，是需要分别追踪的目标；较绿环境中较强的保留关联不能直接换算为植被增加后的定殖收益。
\else
Holding elevation at 100~m, increasing NDVI from 0.30 to 0.76 raised predicted occupancy persistence probability from 29.8\% to 72.0\%, a 42.2-point contrast (95\% interval 9.9--64.9). Colonization changed from 5.3\% to 6.0\%, a 0.7-point contrast with interval $[-2.9,5.3]$. Holding elevation fixed supports an interpretation in which greenness chiefly characterizes retention of occupied sites. Maintaining vegetation conditions at occupied sites and promoting transitions from nonoccupancy to occupancy are therefore distinct monitoring targets; the retention association does not directly quantify a colonization benefit from increasing vegetation.
\fi
\begin{figure}[htbp]\centering
\incfig{fig27_environment_persistence}
\caption{\ifzh
校正模型与等访问次数的报告历史。（a）占域持续概率；（b）定殖；（c）前期有报告后再次报告的比例及计数。前两图标注地点数，阴影表示不足十个地点的组合。过程概率对应模型主要时期之间的转移。
\else
Corrected-model processes and equal-visit reporting histories. Panels show (a) occupancy persistence probability, (b) colonization, and (c) repeated reports conditional on a previous report, with counts. Site counts appear in the first two panels; hatching denotes fewer than ten sites. Process probabilities concern transitions between model primary periods.
\fi}\label{fig:envpersistence}
\end{figure}
\ifzh
等访问次数的记录也显示出环境差异。低海拔低绿度有 1/9 个前期报告地点--时期对再次报告，高绿度为 15/28；高海拔高绿度为 22/32。较绿环境中较高的再次记录比例，与模型估计的较高占域持续概率方向一致；再次记录仍同时受到物种存在与探测机会影响。例如，高海拔高绿度中，前期未报告的十个完整访问对有七个在下一时期报告，而模型定殖概率约为 6\%。较高的后期记录比例与较低的定殖估计并列，说明未报告地点中包含模型判断仍可能已占据的地点。重复调查形成的探测历史有助于区分已占据网格中的再次探测和由未占据到占据的状态转移，从而解释飓风后的报告回升。
\else
Equal-visit histories also differed across environments. Repeated reports occurred in 1/9 previously reporting low-elevation, low-greenness period pairs, 15/28 low-elevation, high-greenness pairs, and 22/32 high-elevation, high-greenness pairs. Higher repeat-report proportions in greener environments are directionally consistent with higher modeled occupancy persistence; repeat reports depend on both presence and detection. Among ten fully sampled high-elevation, high-greenness pairs with no earlier report, seven had a later report, despite modeled colonization near 6\%. Frequent later reports among previously nonreporting pairs alongside low modeled colonization illustrates that nonreporting sites can still be occupied under the detection model. Repeated surveys are therefore valuable for distinguishing renewed detection at occupied grids from transitions out of nonoccupancy.
\fi
\FloatBarrier

